% Generated by publication/build_sections.py; do not edit by hand.
\paragraph{Informal verdict.}
\textbf{A --- Complete and apparently correct.} Confidence: \textbf{high (0.92)}.

\paragraph{Lean coverage.}
Full canonical Q867 coverage. SpaceE defines the exact smooth complex function space with one common derivative bound including n=0 and the exact norm closure F of the span of exp(i lambda x) for |lambda| <= 1. The repaired Fejer interpolation, internal Phillips-Sobczyk obstruction, AP sampling bridge and Tmap reduction unconditionally prove no\_bounded\_projection, F\_not\_complemented and no\_deriv\_invariant\_complement for those concrete spaces.

\subsection{Audited informal answer}
Q867 --- complete solution
\begin{enumerate}
\item Faithful restatement
\end{enumerate}
Let \(E\) be the complex vector space of all \(C^\infty\)-functions \(f:\mathbb R\to\mathbb C\) for which there is a single constant \(M>0\) satisfying

\[
|f^{(n)}(x)|\le M
\qquad
\text{for every }n\in\mathbb N\text{ and }x\in\mathbb R.
\]

Equip \(E\) with

\[
\|f\|_E:=N(f):=\sup_{n\in\mathbb N}\|f^{(n)}\|_\infty .
\]

For \(\lambda\in[-1,1]\), put \(e_\lambda(x)=e^{i\lambda x}\), and let

\[
F=\overline{\operatorname{span}}^{\,\|\cdot\|_E}
   \{e_\lambda:\lambda\in[-1,1]\}.
\]

The question is whether there exists a linear subspace \(G\subset E\) such that

\[
E=F\oplus G,
\]

with \(G\) closed for \(\|\cdot\|_E\) and invariant under differentiation.
This is exactly the statement printed as Q867. Luc Pommeret
The only notational point is that \(0\) must belong to the author's \(\mathbb N\): otherwise constants would have \(N(f)=0\), contradicting the assertion that \(N\) is a norm. Thus \(f^{(0)}=f\) is included.
\begin{enumerate}
\item Answer
\end{enumerate}

Theorem. The subspace \(F\) has no closed topological complement in \(E\).
Consequently, there is in particular no closed complement of \(F\) that is invariant under differentiation.

Thus the answer to Q867 is no. Notice that the conclusion is strictly stronger than requested: differentiation invariance plays no role in the obstruction.
The proof proceeds by transferring any bounded projection \(E\to F\) to a bounded projection

\[
\ell^\infty(\mathbb Z)\longrightarrow AP(\mathbb Z),
\]

where \(AP(\mathbb Z)\) denotes the uniformly almost-periodic sequences. The latter projection cannot exist.

\begin{enumerate}
\item Preliminary Banach-space facts
\end{enumerate}
Lemma 3.1 --- \(E\) is a Banach space
The normed space \((E,\|\cdot\|_E)\) is complete.
Proof
Let \((f_j)\) be Cauchy for \(\|\cdot\|_E\). For each fixed \(n\ge 0\), the sequence \((f_j^{(n)})\) is uniformly Cauchy on \(\mathbb R\), hence converges uniformly to some bounded continuous function \(g_n\).
For \(x,y\in\mathbb R\),

\[
f_j^{(n)}(y)-f_j^{(n)}(x)
   =\int_x^y f_j^{(n+1)}(t)\,dt .
\]

Passing to the limit, using uniform convergence on the compact interval between \(x\) and \(y\), gives

\[
g_n(y)-g_n(x)=\int_x^y g_{n+1}(t)\,dt.
\]

Therefore \(g_n\) is differentiable and \(g_n'=g_{n+1}\). Thus \(f:=g_0\) belongs to \(C^\infty(\mathbb R)\), with

\[
f^{(n)}=g_n\qquad(n\ge 0).
\]

Given \(\varepsilon>0\), choose \(j_0\) so that

\[
\sup_{n\ge0}\|f_j^{(n)}-f_\ell^{(n)}\|_\infty\le\varepsilon
\qquad(j,\ell\ge j_0).
\]

Fixing \(j\ge j_0\) and letting \(\ell\to\infty\), separately for each \(n\), yields

\[
\|f_j^{(n)}-g_n\|_\infty\le\varepsilon
\qquad\text{for every }n.
\]

Hence

\[
\|f_j-f\|_E
 =\sup_n\|f_j^{(n)}-g_n\|_\infty
 \le\varepsilon.
\]

In particular \(f\in E\), and \(f_j\to f\) in \(E\). \(\square\)
The differentiation operator

\[
D:E\to E,\qquad Df=f',
\]

is bounded, since

\[
\|Df\|_E
 =\sup_{n\ge0}\|f^{(n+1)}\|_\infty
 \le \|f\|_E.
\]

Moreover,

\[
De_\lambda=i\lambda e_\lambda,
\]

so \(F\) itself is \(D\)-invariant.
If \(E=F\oplus G\) with \(F,G\) closed, then the projection \(P:E\to F\) along \(G\) is bounded. Indeed, the addition map

\[
F\times G\longrightarrow E,\qquad (f,g)\longmapsto f+g
\]

is a bounded bijection between Banach spaces, so its inverse is bounded by the open mapping theorem.
It therefore suffices to prove:

There is no bounded projection \(P:E\to F\).


\begin{enumerate}
\item A cardinal interpolation operator
\end{enumerate}
Let

\[
h:=8\pi,
\qquad
\psi(w):=
\begin{cases}
\displaystyle\left(\frac{\sin \pi w}{\pi w}\right)^2,&w\ne0,\\[1.2ex]
1,&w=0.
\end{cases}
\]

Then \(\psi\) is entire and

\[
\psi(j-k)=\delta_{jk}
\qquad(j,k\in\mathbb Z).
\]

For \(a=(a_k)_{k\in\mathbb Z}\in\ell^\infty(\mathbb Z)\), define

\[
(Ja)(z):=\sum_{k\in\mathbb Z}a_k\,
     \psi\!\left(\frac zh-k\right).
\tag{4.1}
\]

We shall prove that \(J\) is a bounded map from \(\ell^\infty(\mathbb Z)\) to \(E\).
Lemma 4.1 --- Uniform exponential-type estimate
There is an absolute constant \(C\) such that, for every \(w=u+iv\in\mathbb C\),

\[
\sum_{k\in\mathbb Z}|\psi(w-k)|
 \le C e^{2\pi |v|}.
\tag{4.2}
\]

Proof
First suppose \(|v|\le1\). Choose \(m\in\mathbb Z\) with \(|u-m|\le\frac12\). The function \(\psi(w-m)\) is uniformly bounded on the compact rectangle

\[
|\Re(w-m)|\le\frac12,\qquad |\Im(w-m)|\le1.
\]

For \(k\ne m\),

\[
|\sin\pi(w-k)|=|\sin\pi w|\le e^{\pi|v|}\le e^\pi
\]

and

\[
|w-k|\ge |u-k|\ge |k-m|-\frac12.
\]

Consequently,

\[
\sum_{k\ne m}|\psi(w-k)|
 \le C\sum_{j\ge1}\frac1{(j-\frac12)^2}
 \le C.
\]

Now suppose \(|v|>1\). Since

\[
|\sin\pi w|\le e^{\pi|v|},
\]

we have

\[
\sum_{k\in\mathbb Z}|\psi(w-k)|
 \le \frac{e^{2\pi|v|}}{\pi^2}
      \sum_{k\in\mathbb Z}
      \frac1{(u-k)^2+v^2}.
\]

Uniformly in \(u\),

\[
\sum_{k\in\mathbb Z}
\frac1{(u-k)^2+v^2}
 \le \frac C{|v|}
 \le C,
\]

by comparison with \(\int_{\mathbb R}(t^2+v^2)^{-1}\,dt\). This proves (4.2). \(\square\)
The series in (4.1) converges locally uniformly: on any compact subset, its \(k\)-th term is \(O(k^{-2})\). Hence \(Ja\) is entire. From (4.2),

\[
|(Ja)(x+iy)|
 \le C\|a\|_\infty
       e^{(2\pi/h)|y|}
 =C\|a\|_\infty e^{|y|/4}.
\tag{4.3}
\]

For \(n\ge1\), apply Cauchy's estimate on the circle of radius \(r=4n\) around a real point \(x\). Using (4.3),

\[
\begin{aligned}
|(Ja)^{(n)}(x)|
&\le
\frac{n!}{(4n)^n}
   C\|a\|_\infty e^{n}  \\
&\le
C\|a\|_\infty
   \left(\frac e4\right)^n
\le C\|a\|_\infty,
\end{aligned}
\]

because \(n!\le n^n\) and \(e/4<1\). For \(n=0\), (4.3) on the real axis gives the same conclusion. Therefore

\[
\boxed{\ \|Ja\|_E\le C\|a\|_\infty.\ }
\tag{4.4}
\]

Thus

\[
J:\ell^\infty(\mathbb Z)\longrightarrow E
\]

is bounded.
Define also the sampling map

\[
R:E\longrightarrow\ell^\infty(\mathbb Z),
\qquad
(Rf)_j=f(hj).
\tag{4.5}
\]

It satisfies

\[
\|Rf\|_\infty\le\|f\|_\infty\le\|f\|_E.
\]

By the cardinal property of \(\psi\),

\[
(RJa)_j=(Ja)(hj)
 =\sum_{k\in\mathbb Z}a_k\psi(j-k)
 =a_j.
\]

Hence

\[
\boxed{\ RJ=I_{\ell^\infty(\mathbb Z)}.\ }
\tag{4.6}
\]


\begin{enumerate}
\item Almost-periodic sequences and the interpolation map
\end{enumerate}
Let \(AP(\mathbb Z)\) denote the uniform closure in \(\ell^\infty(\mathbb Z)\) of the trigonometric sequences

\[
k\longmapsto e^{ik\theta},
\qquad \theta\in\mathbb R.
\]

We need two compatibility properties:

\[
R(F)\subset AP(\mathbb Z),
\tag{5.1}
\]

and

\[
J(AP(\mathbb Z))\subset F.
\tag{5.2}
\]

The first is immediate. Indeed,

\[
R(e_\lambda)
   =\bigl(e^{i\lambda hk}\bigr)_{k\in\mathbb Z},
\]

which is a character of \(\mathbb Z\). Since \(R\) is bounded and \(F\) is the \(E\)-closure of the exponential span, (5.1) follows.
The second assertion requires a small entire-function lemma.
Lemma 5.1 --- Entire quasiperiodic functions
Let \(g\) be entire, and suppose that for some \(h>0\), \(\theta\in\mathbb R\) and \(\rho\ge0\),

\[
g(z+h)=e^{i\theta}g(z)
\tag{5.3}
\]

and

\[
|g(x+iy)|\le C e^{\rho|y|}
\qquad(x,y\in\mathbb R).
\tag{5.4}
\]

Then

\[
g(z)=
\sum_{\substack{m\in\mathbb Z\\
|\theta+2\pi m|\le \rho h}}
c_m
\exp\!\left(
i\frac{\theta+2\pi m}{h}z
\right).
\tag{5.5}
\]

In particular, the sum is finite.
Proof
For \(m\in\mathbb Z\), set

\[
\lambda_m:=\frac{\theta+2\pi m}{h}
\]

and define

\[
c_m(y):=
\frac1h\int_0^h
g(x+iy)e^{-i\lambda_m(x+iy)}\,dx.
\]

The integrand

\[
z\longmapsto g(z)e^{-i\lambda_m z}
\]

is \(h\)-periodic by (5.3). Integrating it around a rectangle with horizontal sides at heights \(y_1,y_2\), the vertical integrals cancel. Hence \(c_m(y)\) is independent of \(y\); write it simply as \(c_m\).
If \(\lambda_m>\rho\), take \(y\to-\infty\). Then

\[
|c_m|
 \le C e^{\rho|y|}e^{\lambda_m y}
 =C e^{(\lambda_m-\rho)y}
 \longrightarrow0.
\]

Thus \(c_m=0\). If \(\lambda_m<-\rho\), take \(y\to+\infty\), obtaining

\[
|c_m|
 \le C e^{(\rho+\lambda_m)y}
 \longrightarrow0.
\]

Hence only those \(m\) with \(|\lambda_m|\le\rho\) can contribute.
Subtract the finite sum in (5.5) from \(g\), and multiply by \(e^{-i\theta z/h}\). For every fixed horizontal line this gives a continuous \(h\)-periodic function all of whose Fourier coefficients vanish. Since trigonometric polynomials are dense in the continuous periodic functions, such a function is identically zero. This proves (5.5). \(\square\)
Now fix \(\theta\in\mathbb R\) and put

\[
a_k=e^{ik\theta}.
\]

From the definition of \(J\),

\[
\begin{aligned}
(Ja)(z+h)
 &=\sum_k e^{ik\theta}
      \psi\!\left(\frac zh+1-k\right)\\
 &=e^{i\theta}
   \sum_j e^{ij\theta}
      \psi\!\left(\frac zh-j\right)\\
 &=e^{i\theta}(Ja)(z).
\end{aligned}
\]

Moreover, (4.3) gives

\[
|(Ja)(x+iy)|\le C e^{|y|/4}.
\]

Lemma 5.1, with \(\rho=\frac14\), shows that \(Ja\) is a finite sum of exponentials

\[
e^{i\lambda z}
\qquad\text{with}\qquad
|\lambda|\le\frac14.
\]

Therefore

\[
Ja\in
\operatorname{span}
\{e_\lambda:|\lambda|\le1\}
\subset F.
\]

By linearity, this holds for every trigonometric sequence. Since \(J\) is bounded from the uniform norm to \(\|\cdot\|_E\), passage to the uniform closure gives

\[
\boxed{\ J(AP(\mathbb Z))\subset F.\ }
\tag{5.6}
\]


\begin{enumerate}
\item The essential Banach-space obstruction
\end{enumerate}
Proposition 6.1
There is no bounded projection

\[
\ell^\infty(\mathbb Z)\longrightarrow AP(\mathbb Z).
\]

We give a proof rather than merely invoking the classical general theorem.
6.1. The space \(c_0\) is not complemented in \(\ell^\infty\)
We first work on \(\mathbb N\); relabeling the coordinates gives the same result on \(\mathbb Z\).
Lemma 6.2
Suppose \(T:\ell^\infty\to\ell^\infty\) is bounded and \(T|_{c_0}=0\). Then there is an infinite set \(A\subset\mathbb N\) such that

\[
T|_{\ell^\infty(A)}=0,
\]

where \(\ell^\infty(A)\) denotes the bounded sequences supported on \(A\).
Proof
There exists an uncountable almost-disjoint family \((A_\omega)\) of infinite subsets of a countable set. Explicitly, identify \(\mathbb N\) with the set of finite binary words. For every infinite binary sequence \(\omega\), let \(A_\omega\) be the set of its finite initial segments. Two distinct branches have only finitely many common initial segments.
Suppose that \(T\) is nonzero on every \(\ell^\infty(A_\omega)\). For each \(\omega\), choose \(x_\omega\in\ell^\infty(A_\omega)\) with

\[
\|x_\omega\|_\infty=1,\qquad Tx_\omega\ne0.
\]

Choose integers \(k_\omega,m_\omega\ge1\) such that

\[
|(Tx_\omega)_{k_\omega}|>\frac1{m_\omega}.
\]

Since there are only countably many pairs \((k,m)\), there are fixed \(k,m\) and an uncountable set \(I\) of branches for which

\[
|(Tx_\omega)_k|>\frac1m
\qquad(\omega\in I).
\]

For each \(\omega\in I\), choose a scalar \(\alpha_\omega\) of modulus one so that

\[
\alpha_\omega(Tx_\omega)_k
  =|(Tx_\omega)_k|.
\]

For a finite \(H\subset I\), put

\[
y=\sum_{\omega\in H}\alpha_\omega x_\omega.
\]

The union of the pairwise intersections

\[
B=\bigcup_{\substack{\omega,\eta\in H\\\omega\ne\eta}}
(A_\omega\cap A_\eta)
\]

is finite. Write \(y=u+v\), where \(v\) is the restriction of \(y\) to \(B\). Then \(v\) is finitely supported, so \(Tv=0\), while outside \(B\) at most one of the \(x_\omega\) is nonzero. Hence

\[
\|u\|_\infty\le1.
\]

It follows that

\[
\|Ty\|_\infty=\|Tu\|_\infty\le\|T\|.
\]

On the other hand,

\[
(Ty)_k
 =\sum_{\omega\in H}|(Tx_\omega)_k|
 >\frac{|H|}{m}.
\]

Taking \(|H|>m\|T\|\) gives a contradiction. \(\square\)
If \(P:\ell^\infty\to c_0\) were a bounded projection, then \(T=I-P\) would vanish on \(c_0\). Lemma 6.2 would give an infinite \(A\) such that \(P\) is the identity on \(\ell^\infty(A)\). In particular,

\[
P1_A=1_A.
\]

But \(P1_A\in c_0\), whereas \(1_A\notin c_0\). This is impossible.
Thus

\[
\boxed{\ c_0\text{ is not complemented in }\ell^\infty.\ }
\tag{6.1}
\]

This is the classical Phillips--Sobczyk theorem; the almost-disjoint-family proof above is also documented in the cited lecture notes. IM PAN
6.2. Injectivity of \(\ell^\infty\)
Recall that a Banach space \(Y\) is called injective if every bounded operator \(A:X_0\to Y\), defined on a closed subspace \(X_0\) of a Banach space \(X\), extends to a bounded operator \(X\to Y\).
The space \(\ell^\infty\) is injective. Indeed, write

\[
A(x)=(a_n(x))_{n\ge1},
\qquad a_n\in X_0^*,
\qquad \|a_n\|\le\|A\|.
\]

By Hahn--Banach, extend every \(a_n\) to \(\widetilde a_n\in X^*\) with the same norm, and set

\[
\widetilde A(x)
   =(\widetilde a_n(x))_{n\ge1}.
\]

Then \(\widetilde A:X\to\ell^\infty\) is a bounded extension of \(A\).
A complemented subspace of an injective Banach space is injective: extend into the ambient injective space, then compose with the bounded projection.
Consequently, no complemented subspace of \(\ell^\infty\) can be isomorphic to a Banach space that has a complemented copy of \(c_0\). Otherwise that copy of \(c_0\) would be injective; applying injectivity to the identity map

\[
c_0\subset\ell^\infty\longrightarrow c_0
\]

would produce a projection \(\ell^\infty\to c_0\), contradicting (6.1).
6.3. A complemented copy of \(c_0\) in \(C(\mathbb T)\)
Choose distinct points \(x_n\in\mathbb T\) converging to \(x_0\), and pairwise disjoint open arcs \(U_n\) such that

\[
x_n\in U_n,\qquad x_0\notin\overline{U_n},
\]

and the only accumulation point of the \(U_n\)'s is \(x_0\).
Choose continuous functions \(\varphi_n:\mathbb T\to[0,1]\) satisfying

\[
\operatorname{supp}\varphi_n\subset U_n,
\qquad
\varphi_n(x_n)=1.
\]

Define

\[
A:c_0\to C(\mathbb T),
\qquad
A(c)=\sum_{n\ge1}c_n\varphi_n.
\]

Because the supports are disjoint,

\[
\left\|
\sum_{n>N}c_n\varphi_n
\right\|_\infty
\le\sup_{n>N}|c_n|\longrightarrow0.
\]

Thus the series converges uniformly and

\[
\|A(c)\|_\infty=\|c\|_\infty.
\]

Define

\[
B:C(\mathbb T)\to c_0,
\qquad
Bf=(f(x_n)-f(x_0))_{n\ge1}.
\]

Continuity of \(f\) and \(x_n\to x_0\) imply \(Bf\in c_0\), and

\[
\|Bf\|_\infty\le2\|f\|_\infty.
\]

Moreover,

\[
BA=I_{c_0}.
\]

Hence \(AB\) is a bounded projection of \(C(\mathbb T)\) onto the copy \(A(c_0)\).
It follows from the preceding injectivity argument that:

\[
\boxed{\text{No complemented subspace of }\ell^\infty
       \text{ is isomorphic to }C(\mathbb T).}
\tag{6.2}
\]

6.4. A complemented copy of \(C(\mathbb T)\) in \(AP(\mathbb Z)\)
We construct the relevant part of the Bohr compactification explicitly.
For \(n\in\mathbb Z\), define

\[
\iota(n)
   =(\zeta^n)_{\zeta\in\mathbb T}
   \in\mathbb T^{\mathbb T},
\]

and let

\[
b\mathbb Z
   :=\overline{\iota(\mathbb Z)}
   \subset\mathbb T^{\mathbb T}.
\]

By Tychonoff's theorem, \(b\mathbb Z\) is a compact Hausdorff group.
For every \(\zeta\in\mathbb T\), let

\[
p_\zeta(x)=x_\zeta,
\qquad x\in b\mathbb Z.
\]

The finite linear span of the \(p_\zeta\)'s is a self-adjoint unital algebra:

\[
p_\zeta p_\eta=p_{\zeta\eta},
\qquad
\overline{p_\zeta}=p_{\bar\zeta},
\qquad
p_1=1.
\]

It separates the points of \(b\mathbb Z\), so by the complex Stone--Weierstrass theorem it is uniformly dense in \(C(b\mathbb Z)\).
Since \(\iota(\mathbb Z)\) is dense, restriction gives an isometric map

\[
C(b\mathbb Z)\longrightarrow \ell^\infty(\mathbb Z),
\qquad
\Phi\longmapsto
\bigl(\Phi(\iota(n))\bigr)_{n\in\mathbb Z}.
\]

The coordinate function \(p_\zeta\) restricts to \(n\mapsto\zeta^n\). It follows that the image is exactly \(AP(\mathbb Z)\). Thus

\[
AP(\mathbb Z)\cong C(b\mathbb Z).
\tag{6.3}
\]

Now choose \(\alpha\in\mathbb R\) with \(\alpha/(2\pi)\) irrational. The coordinate map

\[
\pi:b\mathbb Z\to\mathbb T,
\qquad
\pi(x)=x_{e^{i\alpha}},
\]

is a continuous group homomorphism. Its image contains

\[
\{e^{in\alpha}:n\in\mathbb Z\},
\]

which is dense in \(\mathbb T\). Since the image of a compact space is compact and hence closed, \(\pi\) is surjective.
Let \(K=\ker\pi\), and let \(m_K\) be normalized Haar measure on the compact group \(K\). Define

\[
\mathcal E\Phi(x)
   :=\int_K\Phi(xk)\,dm_K(k).
\]

Then \(\mathcal E\) is a norm-one projection from \(C(b\mathbb Z)\) onto the \(K\)-invariant functions. These functions are precisely the functions factoring through the quotient

\[
b\mathbb Z/K\cong\mathbb T.
\]

Thus \(\pi^*C(\mathbb T)\) is complemented in \(C(b\mathbb Z)\). Under (6.3), it corresponds to

\[
H_\alpha
 =\left\{
 \bigl(f(e^{in\alpha})\bigr)_{n\in\mathbb Z}:
 f\in C(\mathbb T)
 \right\}.
\]

The map

\[
C(\mathbb T)\to H_\alpha,
\qquad
f\longmapsto(f(e^{in\alpha}))_{n\in\mathbb Z},
\]

is isometric because the orbit \(\{e^{in\alpha}\}\) is dense.
Therefore \(AP(\mathbb Z)\) contains a complemented isometric copy of \(C(\mathbb T)\).
If \(AP(\mathbb Z)\) were complemented in \(\ell^\infty(\mathbb Z)\), this copy of \(C(\mathbb T)\) would also be complemented in \(\ell^\infty(\mathbb Z)\), contradicting (6.2). This proves Proposition 6.1. \(\square\)
For comparison, the general theorem for discrete groups says that \(AP(G)\) is complemented in \(\ell^\infty(G)\) exactly when \(AP(G)\) is finite-dimensional; Lau and Losert explicitly record this result of Takahashi. MSP

\begin{enumerate}
\item Completion of the proof
\end{enumerate}
Assume, toward a contradiction, that \(F\) has a closed complement in \(E\). Let

\[
P:E\longrightarrow F
\]

be the associated bounded projection.
Define

\[
Q:=RPJ:
\ell^\infty(\mathbb Z)\longrightarrow\ell^\infty(\mathbb Z).
\]

For arbitrary \(a\in\ell^\infty(\mathbb Z)\), \(PJa\in F\), and therefore, by (5.1),

\[
Qa=R(PJa)\in AP(\mathbb Z).
\]

If \(a\in AP(\mathbb Z)\), then by (5.6),

\[
Ja\in F.
\]

Hence \(PJa=Ja\), and by (4.6),

\[
Qa=RPJa=RJa=a.
\]

Thus \(Q\) is a bounded projection of \(\ell^\infty(\mathbb Z)\) onto \(AP(\mathbb Z)\). This contradicts Proposition 6.1.
Therefore no bounded projection \(E\to F\) exists, and \(F\) has no closed complement in \(E\). In particular it has no closed differentiation-invariant complement. \(\square\)

\begin{enumerate}
\item Boundary checks and the role of the hypotheses
\end{enumerate}
For \(|\lambda|\le1\),

\[
e_\lambda^{(n)}=(i\lambda)^n e_\lambda,
\]

so

\[
\|e_\lambda\|_E=1.
\]

The endpoints \(\lambda=\pm1\) are admissible: every derivative still has supremum norm \(1\). Conversely, if \(|\lambda|>1\), then

\[
\|e_\lambda^{(n)}\|_\infty=|\lambda|^n\to\infty,
\]

so \(e_\lambda\notin E\). Thus the spectral interval appearing in the problem is exactly compatible with the common derivative bound.
The subspace \(F\) is proper. Indeed, consider

\[
s(x)=\int_{-1/2}^{1/2}e^{itx}\,dt
=
\begin{cases}
\dfrac{\sin(x/2)}{x/2},&x\ne0,\\[1ex]
1,&x=0.
\end{cases}
\]

Then

\[
s^{(n)}(x)
 =\int_{-1/2}^{1/2}(it)^n e^{itx}\,dt,
\]

and consequently

\[
\|s^{(n)}\|_\infty
 \le \int_{-1/2}^{1/2}|t|^n\,dt
 =\frac1{2^n(n+1)}
 \le1.
\]

Thus \(s\in E\). But \(s(x)\to0\) as \(|x|\to\infty\), while a nonzero uniformly almost-periodic function cannot vanish at infinity: its sufficiently small almost-periods are relatively dense, so a nonzero value recurs arbitrarily far away. Since \(F\) is contained in the uniformly almost-periodic functions on \(\mathbb R\), \(s\notin F\).
The conditions enter as follows.


Differentiation invariance of the proposed complement is unnecessary. The proof rules out every closed complement.


Closedness is decisive. It produces a bounded projection. Without closedness, every vector subspace has an algebraic complement in ZFC, although such an algebraic complement need not be differentiation-invariant.


The proof uses only the smaller frequency band \([-\frac14,\frac14]\subset[-1,1]\). Thus the obstruction comes from the presence of a genuine interval of frequencies, not from the endpoints \(\pm1\).


The single uniform bound over all derivative orders is used in the norm estimate for the interpolation operator \(J\).


No computation is involved.

\begin{enumerate}
\item Topological and set-theoretic assumptions
\end{enumerate}
All normed spaces are complex, Hausdorff and equipped with their norm topologies. Their closures and convergence arguments are metrizable, so sequences suffice throughout the Banach-space portions of the proof.
The compact group \(b\mathbb Z\), however, is generally nonmetrizable. Its construction uses the full product \(\mathbb T^{\mathbb T}\), Tychonoff compactness, the complex Stone--Weierstrass theorem, and Haar measure on a compact Hausdorff group. No appeal to sequential compactness of \(b\mathbb Z\) is made.
The countable structures are the sampling set \(\mathbb Z\), the sequence spaces \(\ell^\infty(\mathbb Z)\) and \(c_0(\mathbb Z)\), and the convergent sequence used inside \(C(\mathbb T)\). The almost-disjoint family is uncountable but consists of subsets of a countable binary tree.
The proof is carried out in ordinary ZFC. The standard choice-dependent ingredients are Tychonoff's theorem, Hahn--Banach, Haar measure, the open mapping theorem, and the selections made in the almost-disjoint-family argument. No additional hypothesis such as the continuum hypothesis is used.

\[
\boxed{\text{Q867 has a negative answer: }F\text{ is not complemented in }E.}
\]
